\chapter{Introduction}
\label{chap:introduction}

\section{Comparing dependence across populations}

Many empirical studies ask whether the relationship between two variables
changes across populations. A treatment may alter the association between a
biomarker and an outcome; environmental conditions may change the dependence
between two species; or two periods may exhibit different coupling between
components of a complex system. In each case, the quantity of interest is the
strength of association, not only the marginal behaviour of either variable.

Mutual information (MI) provides a general measure of statistical dependence
\citep{shannon1948}. For categorical variables, it compares the observed joint
distribution with the joint distribution that would arise under independence.
Unlike correlation, MI does not restrict dependence to a linear or monotone
form. This generality makes MI useful in complex-systems research and in
applications where the structure of an association is not known in advance
\citep{brillinger2004,lizier2014}.

The most familiar significance test for categorical MI asks whether two
variables are independent within one population. The corresponding null is
\(
I(X;Y)=0
\), and the likelihood-ratio statistic has an asymptotic chi-squared
distribution under standard fixed-alphabet conditions. A permutation test can
address the same question by shuffling one variable and thereby destroying its
association with the other.

Comparing dependence across two populations is a different inferential
problem. Let \(P\) and \(Q\) denote the joint distributions of the same aligned
categorical variables in two independent populations. This thesis studies
\begin{equation}
  H_0:I(P)=I(Q)
  \quad\text{against}\quad
  H_1:I(P)\ne I(Q).
  \label{eq:intro-null}
\end{equation}
The complete distributions may differ under this null. Their margins, cell
probabilities, and patterns of association need not be equal. Equation
\eqref{eq:intro-null} therefore defines a weak null about one scalar functional
of each distribution.

This distinction determines the appropriate analysis. Applying a separate
chi-squared independence test to each table does not test whether the two MI
values are equal. Shuffling one variable within either population also targets
independence rather than Equation~\eqref{eq:intro-null}. The required method
must estimate the signed difference \(I(P)-I(Q)\) and quantify its sampling
uncertainty.

\section{The finite-sample problem}

For an \(r\times c\) contingency table, the plug-in MI estimator replaces each
population probability by its observed proportion. The estimator is simple
and consistent, but it is nonlinear and has positive finite-sample bias
\citep{miller1955,moddemeijer1989,paninski2003}. The leading bias depends on
the alphabet dimensions and sample size. Consequently, two groups with equal
population MI can have differently biased estimates when their sample sizes
differ.

Away from independence, MI is a smooth functional of the cell probabilities.
A first-order Taylor expansion expresses the estimation error as an average of
centred pointwise-MI contributions. This representation gives an analytic
standard error and leads to a normal Wald test for Equation
\eqref{eq:intro-null}. The calculation is deterministic and requires one pass
over each table.

The normal Wald test treats the estimated standard error as sufficiently
stable. That approximation becomes less reliable when the sample is small,
the margins are skewed, or many cells have low expected counts. In these
settings, repeated samples can produce substantially different estimates of
the MI variance. Ignoring that second layer of uncertainty can make the
normal-reference test reject too often.

Welch's unequal-variance test addresses an analogous issue when comparing two
means. It forms a Student statistic and uses effective degrees of freedom that
depend on the two estimated variance contributions
\citep{satterthwaite1946,welch1947}. Hutcheson transported this architecture
to the comparison of Shannon diversity values \citep{hutcheson1970}. A direct
transport to MI, however, does not fully describe the uncertainty of the MI
variance estimate. Changing one observation alters the joint probabilities,
the margins, every affected pointwise-MI value, and the estimated MI variance.

\section{Aim and research questions}

This thesis develops and evaluates an MI-specific Welch--Satterthwaite
reference for the weak null in Equation~\eqref{eq:intro-null}. The proposed
test keeps the bias-corrected MI difference and first-order standard error of
the normal Wald test. It changes only the reference distribution used to
calibrate the standardized statistic.

Figure~\ref{fig:intro-method-overview} gives the high-level calculation. The
two tables are processed separately until their bias-corrected MI estimates
and variance contributions are combined. The proposed contribution lies in
the final Welch-type calibration: its degrees of freedom are derived from the
sampling stability of the complete MI variance estimator.

\begin{figure}[htbp]
  \centering
  \includegraphics[width=0.96\textwidth]{method_overview.pdf}
  \caption{High-level analytic pipeline for comparing MI between two
  independent count tables. Expanded Welch uses the same effect estimate and
  standard error as the two baselines but supplies MI-specific component
  degrees of freedom for the final Student reference.}
  \label{fig:intro-method-overview}
\end{figure}

The study addresses three research questions.

\begin{enumerate}[leftmargin=*,label=\textbf{RQ\arabic*:}]
  \item \textbf{Method construction.}
  How can the Welch--Satterthwaite principle be adapted to the nonlinear
  plug-in MI variance estimator?

  \item \textbf{Statistical performance.}
  Does the expanded method improve type-I error calibration and confidence
  interval coverage relative to normal Wald and simple
  Welch--Satterthwaite inference across well-sampled, skewed, sparse, and
  unequal-sample regimes?

  \item \textbf{Practical trade-offs.}
  What power, computational cost, validity-rate, and operating-range trade-offs
  accompany the correction?
\end{enumerate}

\section{Contributions}

The thesis makes four contributions.

First, it derives component degrees of freedom for the estimated first-order
MI variance. The derivation treats that variance as a complete functional of
the joint distribution, including the dependence of pointwise MI on the row
and column margins.

Second, it combines these component degrees of freedom into a deterministic
expanded Welch--Satterthwaite test. The complete calculation requires
\(O(rc)\) operations for an \(r\times c\) table and does not require
resampling.

Third, it evaluates normal Wald, simple Welch--Satterthwaite, and expanded
Welch--Satterthwaite inference on a common grid of exactly equal-MI population
pairs. The grid separates well-sampled controls from skewed, rare-cell,
ultra-sparse, and widespread-sparsity regimes.

Fourth, it characterises the method's operating range. The results identify
regimes in which the correction improves calibration, quantify its power and
runtime costs, and show where first-order inference remains unreliable.

The contribution is deliberately bounded. Welch--Satterthwaite moment
matching, influence-function theory, bias correction for plug-in information
measures, and comparisons of MI-related indices all have established
literatures. The new element is the MI-specific variance-influence
degrees-of-freedom construction and its controlled finite-sample evaluation
for the equal-MI weak null.

\section{Scope}

The analysis assumes two independent samples of aligned categorical pairs.
The principal theory concerns fixed finite alphabets, positive population cell
probabilities, and positive first-order MI variances. The target MI values in
the primary validation are positive and bounded away from zero.

Testing independence is outside the method's regular operating range. At
\(I(P)=0\), the first-order MI influence function is zero and the usual
normal expansion degenerates. The chi-squared independence limit then arises
from a second-order expansion. Appendix~\ref{app:independence-boundary} derives this
boundary explicitly.

The thesis does not develop paired-sample inference, conditional MI, transfer
entropy, or a universal sparse-table correction. These extensions require
additional dependence structures or different asymptotic arguments.

\section{Thesis structure}

Chapter~\ref{chap:background} introduces the information-theoretic and
statistical foundations, reviews prior comparison methods, and identifies the
specific methodological gap. Chapter~\ref{chap:baselines} defines the
two-sample problem and derives the shared bias-corrected statistic, normal Wald
baseline, and simple Welch baseline. Chapter~\ref{chap:expanded} derives the
expanded MI-specific degrees of freedom.

Chapter~\ref{chap:design} presents the population construction, simulation
grid, evaluation metrics, and reproducibility controls. Chapter
\ref{chap:results} reports calibration, coverage, power, validity, and runtime.
Chapter~\ref{chap:discussion} interprets the findings, defines the operating
range, and discusses limitations. Chapter~\ref{chap:conclusion} answers the
research questions and concludes the thesis.
